% 06-conclusion.tex
\chapter{Fazit und Ausblick}
\label{chap:conclusion}

Dieses Kapitel schließt die Arbeit ab. Es beinhaltet eine Zusammenfassung der wichtigsten Ergebnisse, eine ehrliche Betrachtung der Einschränkungen (Limitationen) sowie Ideen für Folgearbeiten.

\section{Zusammenfassung}
In dieser Arbeit wurde gezeigt, wie \dots (Hier werden die Forschungsfragen aus der Einleitung wieder aufgegriffen und beantwortet). Die entwickelten Methoden haben bestätigt, dass sich der untersuchte Prozess signifikant optimieren lässt.

\section{Limitationen}
Trotz der vielversprechenden Ergebnisse unterliegt die vorliegende Arbeit gewissen Einschränkungen, die bei der Interpretation der Resultate berücksichtigt werden müssen:
\begin{itemize}
  \item \textbf{Datengrundlage:} Die Evaluierung fand auf einem verhältnismäßig kleinen Datensatz statt, was die Generalisierbarkeit der Ergebnisse einschränken könnte.
  \item \textbf{Hardware-Restriktionen:} Aufgrund limitierter Rechenleistung konnten bestimmte komplexe Modelle nicht vollumfänglich trainiert werden.
\end{itemize}

\section{Ausblick}
Aufbauend auf den Erkenntnissen dieser Arbeit bieten sich diverse Anknüpfungspunkte für zukünftige Forschungen. Zukünftig könnte man beispielsweise untersuchen, wie sich das System bei der Integration von Echtzeit-Datenströmen verhält.

\clearpage